\section{Supplement to Section 6: basin-aware gain scheduling}\label{supp:ba}

This section gives the protocol of the basin-aware scan of Section~\ref{M-sec:basin-aware} and its secondary results.

\subsection{Protocol of the scan}\label{supp:ba-protocol}

\paragraph{Grid} The depths are $\mathcal A=\{2,4,\dots,32\}$\,\si{\milli\metre}; the speeds are 500--\SI{3000}{\rpm} in steps of \SI{125}{\rpm} (21 speeds) and the speeds 1090, 1420 and \SI{2870}{\rpm} of Section~\ref{M-sec:bistab-basin}. The gains are first those of $\{0,250,\dots,3000\}$\,\si{\volt\second\per\metre}, then, for each $\beta_{\mathrm{req}}$, gains $\pm125$ and then $\pm62.5$\,\si{\volt\second\per\metre} around the best gain so far.

\paragraph{Conditions} Linear stability uses the zero-order limit $\alim(n,g)$. The voltage margin uses the rigorous bound $\bar b$ of \ref{M-app:lmi} (the largest of its 16 segment bounds, harmonics above the 40th included), $\bar b\le0.9\Vmax=\SI{8.1}{\volt}$; between grid depths it was also checked on a \SI{0.25}{\milli\metre} grid. The certified kick is that of C2 with $e=e_2$, the validity bound $\delta=0.5$ and the \SI{9}{\volt} clip, rescaled to $\bar b$ as in Section~\ref{M-sec:cert-lmi}. The binding constraint reported for a speed is the one that fails at the next grid depth, $a_{\mathrm{rob}}+\SI{2}{\milli\metre}$.

\paragraph{Cost} For each gain the depths are taken upwards from \SI{2}{\milli\metre} until one is not admissible for $\beta_{\mathrm{req}}=\SI{1}{\milli\metre\per\second}$. The closed-form conditions bound $a_{\mathrm{rob}}(n,g)$ from above before any certificate is computed, so gains that cannot improve on the best depth found are skipped (branch and bound). The scan uses the collocation order $N_{\mathrm{red}}=N-8$, whose linear limit is within $1.2\times10^{-4}$ of the zero-order limit, and the solver defaults. Every selected design point, every lower grid depth whose reduced-order kick is below $1.15\,\beta_{\mathrm{req}}$ and the failing depth were then certified again with the default order $N$ (Section~\ref{M-sec:cert-model}), with the default and, where needed, the strict solver settings, the larger certificate being kept; $a_{\mathrm{rob}}$ is lowered where a default-order certificate falls below $\beta_{\mathrm{req}}$.

\paragraph{Validation} At 750, 1090, 1420, 2000, 2500 and \SI{2870}{\rpm}, the simulated critical kick of the design point (face entry, full surface memory, \SI{9}{\volt} clip, kick at revolution 100, classifier window of 300 revolutions) was computed at four phases, 0, 4, 8 and 12 of 16, and kicks of exactly $\beta_{\mathrm{req}}$ were applied at all 16 phases; the same critical kick was computed at the checked depth $\achk$ of the authority-only design for comparison.
% src: verification/paper/basin_aware.py (docstring: DEPTHS, GRID_RPM, C1_RPM, G_COARSE, REFINE, MARGIN 1.15, N_red = nodes_for - 8, c2.FAST/TIGHT, VAL_SPEEDS, VAL_PHASES; recert); verification/paper/basin_aware_nconv.json (reduced order within 1.2e-4 of the ZOA limit)

\NEWBASUPP
